\chapter{版本记录与项目分工}

\section{版本更新记录}

本节记录语界 LinguaSpace 多语文旅智能导览与人才培养系统在文档编制、需求分析、系统设计、测试说明和用户手册完善过程中的主要版本变化。版本号采用“主版本.次版本”形式，其中主版本表示阶段性文档交付，次版本表示同一阶段内的内容补充、结构调整或排版修正。

\begin{longtable}{@{}p{1.7cm}p{2.5cm}p{3.0cm}p{5.6cm}@{}}
\toprule
版本号 & 更新时间 & 更新类型 & 更新内容 \\
\midrule
V0.1 & - & 初始整理 & 完成项目资料梳理，确定系统名称、目标用户、核心业务范围和文档总体结构。 \\
V0.2 & - & 需求完善 & 补充游客端、学生端、导游端和管理端的功能需求，明确非功能需求、数据需求和业务流程。 \\
V0.3 & - & 图表补充 & 根据需求内容绘制用例图、业务流程图、ER 图、类图和顺序图，并将图表插入需求章节。 \\
V0.4 & - & 系统设计 & 完成概要设计和详细设计，补充分层架构、包结构、接口表、核心领域类图和数据库结构说明。 \\
V0.5 & - & 架构优化 & 优化系统架构图、服务关系图和包依赖图，使图示更符合软件工程文档表达习惯。 \\
V0.6 & - & 测试文档 & 完善单元测试、集成测试、接口测试、系统测试、异常边界测试、非功能测试和测试结果记录。 \\
V0.7 & - & 用户手册 & 完善软件概述、运行环境、客户端进入方式、功能操作步骤和界面截图占位说明。 \\
V0.8 & - & 交付整理 & 补充版本更新记录和组员分工表，统一文档结构、表格表达和 PDF 编译结果。 \\
V1.0 & - & 正式交付 & 形成完整的软件工程综合报告，包含需求、设计、测试、用户使用说明和项目记录。 \\
\bottomrule
\end{longtable}

\section{组员分工表}

项目分工按照软件工程文档编制和系统功能模块划分。每名成员负责对应模块的资料整理、功能分析、图表绘制、文档撰写和交叉检查。实际提交前可将“组员一”等占位替换为真实姓名。

\begin{longtable}{@{}p{2.0cm}p{2.5cm}p{4.2cm}p{4.8cm}@{}}
\toprule
成员 & 角色 & 主要负责内容 & 交付成果 \\
\midrule
组员一 & 项目统筹 & 负责总体选题、系统目标梳理、文档结构规划、章节衔接和最终 PDF 汇总。 & 项目总体说明、目录结构、文档整合、最终编译检查。 \\
组员二 & 需求分析 & 负责用户角色、功能需求、非功能需求、用例模型、业务流程和数据需求整理。 & 需求分析章节、用例说明、业务流程图、需求追踪表。 \\
组员三 & 系统设计 & 负责概要设计、详细设计、分层架构、包结构、核心类设计、接口设计和数据库结构设计。 & 系统设计章节、架构图、类图、接口表、数据库表结构。 \\
组员四 & 测试设计 & 负责测试范围、测试环境、单元测试、集成测试、接口测试、系统测试和测试结果汇总。 & 软件测试文档、测试用例表、测试结果表、缺陷统计表。 \\
组员五 & 用户手册 & 负责软件概述、运行环境、客户端入口、功能使用步骤、常见问题和截图占位整理。 & 用户使用说明书、操作步骤说明、界面截图预留说明。 \\
组员六 & 图表与排版 & 负责 PlantUML 图表整理、图片命名、图表插入、版式检查、Overfull 问题修正和 PDF 编译。 & UML 图、架构图、流程图、排版修正记录、编译后的 PDF。 \\
\bottomrule
\end{longtable}

\section{协作与审核记录}

\begin{longtable}{@{}p{3.0cm}p{4.2cm}p{6.0cm}@{}}
\toprule
协作事项 & 参与人员 & 处理内容 \\
\midrule
需求评审 & 全体组员 & 对系统目标、角色边界、功能范围和非功能约束进行统一确认。 \\
设计评审 & 系统设计负责人、需求分析负责人、图表与排版负责人 & 检查架构设计、类图、接口表和数据库结构是否覆盖需求内容。 \\
测试评审 & 测试设计负责人、系统设计负责人 & 检查测试项是否覆盖核心功能链路、异常场景、接口契约和非功能要求。 \\
文档统稿 & 项目统筹负责人、用户手册负责人、图表与排版负责人 & 统一章节命名、表格风格、图表编号、截图占位和最终 PDF 输出格式。 \\
最终检查 & 全体组员 & 检查内容完整性、格式一致性、图表可读性和编译结果。 \\
\bottomrule
\end{longtable}
